\chapter{The Archimedean Wilson-Surface Operator \(F_{\infty}\)}
\label{v4-arith:w9-r2:archimedean-wilson-surface}

\emph{Chapter~\ref{v4-arith:w8-a:deninger-chiral-homology} isolated
the conditional archimedean component \(F_{\infty}\) of the formal Frobenius
\(F=F_{\infty}\otimes\bigotimes'_{p}\psi_{p}\) in three guises:
(i) the Arakelov scaling exponential
\(F_{\infty}=\exp(-\Theta_{\infty})\), where
\(\Theta_{\infty}=\tfrac12\mathrm{id}+i\widehat D\) and
\(\widehat D=-i\partial_t\) is the self-adjoint multiplicative
scaling generator of \Cref{v4-arith:w9-r2:rem:t-coord};
(ii) the bulk area-scaling Wilson-surface operator at \(\infty\) inside
the conjectural 3D arithmetic Chern-Simons of
Chapter~\ref{v4-arith:3d-acs-E1bar}; (iii) the Mellin-conjugate of the
arithmetic Virasoro zero-mode \(L_{0}^{\mathrm{Ar}}\) under the
archimedean Zhu correspondence of Chapter~\ref{v4-arith:rh-reformulation}.
The construction of (i) is elementary and
unconditional. The realisation (ii) is a quantum-ACS observable
supplied conditionally on the level-\(k^{\mathrm{Ar}}\) bulk measure of
\Cref{v4-arith:3d-acs:conj:qu-acs}. The conjugacy (iii) with
\(L_{0}^{\mathrm{Ar}}\) is a theorem on the archimedean Fock fibre
after the stated primary-Virasoro comparison.
The determinant-divisor identification then follows: under
\Cref{v4-arith:rh:hyp:HPAr}, the Tate-shifted Mellin image of
\(L_{0}^{\mathrm{Ar}}\) on \(H^{1}_{\mathrm{ch}}\) has divisor
\(\operatorname{div}\xi_{\mathbb{R}}\), while the determinant divisor
of \(F_{\infty}\) is its exponential pushforward. Archimedean purity for
\(F_{\infty}\) on \(H^{1}_{\mathrm{ch}}\) is equivalent to the A-TP
positivity of Chapter~\ref{v4-arith:w5-a:chapter}, and the latter is
equivalent to RH.}

\section{Arakelov Archimedean Structure}
\label{v4-arith:w9-r2:arakelov-structure}

The archimedean fibre of \(\overline{\mathrm{Spec}\,\mathbb{Z}}\) is the
compactifying place at \(\infty\), where \(\mathbb{Z}\hookrightarrow\mathbb{R}\)
is replaced by the Arakelov datum: the archimedean completion
\(\mathbb{Q}_{\infty}=\mathbb{R}\), the normalised absolute value
\(|x|_{\infty}=|x|\), and a choice of smooth Green metric on
\(\mathbb{P}^{1}_{\mathbb{C}}=\mathbb{C}\cup\{\infty\}\) when the
arithmetic surface is projective. The archimedean fibre relevant to
the arithmetic bar-character programme is the complex plane
\(\mathbb{C}_{\infty}=\mathbb{C}\), equipped with the standard Lebesgue
Haar measure \(d^{2}x=dx\,dy\) on the additive side and the scale-invariant
Haar measure \(d^{\times}\lambda=d\lambda/\lambda\) on
\(\mathbb{R}^{\times}_{+}\) on the multiplicative side.

\begin{definition}[Arakelov scaling action]
\label{v4-arith:w9-r2:def:scaling}
\ClaimStatusDefinitional. Set \(\mathbb{R}^{\times}_{+}\) with coordinate
\(\lambda\). The \emph{Arakelov scaling action} on the archimedean fibre
is
\begin{equation}
\label{v4-arith:w9-r2:eq:scaling}
\rho_{\lambda}\;\colon\;(\mathbb{C}_{\infty},|\,\cdot\,|_{\infty})
\longrightarrow (\mathbb{C}_{\infty},|\,\cdot\,|_{\infty}),
\qquad
x\longmapsto \lambda\cdot x,
\end{equation}
together with the induced action on the multiplicative Haar measure
\(\rho_{\lambda}^{*}(d^{\times}|x|)=d^{\times}|x|\). On
\(L^{2}(\mathbb{R}^{\times}_{+},d^{\times}x)\) the unitary scaling
action is
\[
 (U_{\lambda}f)(x):=f(\lambda^{-1}x).
\]
The self-adjoint scaling generator is
\[
 \widehat D:=-i\frac{d}{dt}\bigg|_{t=0}U_{e^t}.
\]
The half-density factor \(|\lambda|^{-1/2}\) belongs to the additive
Lebesgue model \(L^{2}(\mathbb{R}_{+},dx)\), not to the multiplicative
Haar model.
\end{definition}

\begin{remark}[\(t\)-coordinate]
\label{v4-arith:w9-r2:rem:t-coord}
Put \(t=\log\lambda\) so that \(\lambda=e^{t}\) and \(d^{\times}\lambda=dt\);
then the scaling action becomes translation
\((U_{e^t}f)(u)=f(u-t)\) on \(L^{2}(\mathbb{R},du)\) under
\(u=\log|x|\), and the generator is
\[
 \widehat D=-i\partial_u .
\]
The Tate half-weight is not part of the unitary translation
generator. It enters through the normal operator
\[
 \Theta_{\infty}:=\frac12\mathrm{id}+i\widehat D,
\]
whose Mellin eigenvalue is \(s=\tfrac12+i\tau\).
\end{remark}

\begin{lemma}[Mellin unitarity]
\label{v4-arith:w9-r2:lem:mellin}
\ClaimStatusProvedHere. The Mellin transform
\(\mathcal{M}f(\tau):=\int_{0}^{\infty}f(x)\,x^{-i\tau}\,d^{\times}x\)
is a unitary isomorphism
\begin{equation}
\label{v4-arith:w9-r2:eq:mellin}
\mathcal{M}\;\colon\; L^{2}(\mathbb{R}^{\times}_{+},d^{\times}x)
\xrightarrow{\ \ \sim\ \ } L^{2}\bigl(\mathbb{R}_{\tau},\tfrac{d\tau}{2\pi}\bigr),
\end{equation}
conjugating the multiplicative-scaling generator \(\widehat{D}\) to
multiplication by \(\tau\). After the Tate shift
\(s=\tfrac12+i\tau\), the normal operator
\(\Theta_{\infty}=\tfrac12\mathrm{id}+i\widehat D\) is multiplication
by \(s\) on the critical line.
\end{lemma}

\begin{proof}
Standard: Plancherel for the multiplicative group
\(\mathbb{R}^{\times}_{+}\simeq\mathbb{R}\) via \(t=\log x\). The additive
Fourier transform on \(L^{2}(\mathbb{R},dt)\) conjugates
\(-i\partial_{t}\) to multiplication by the real frequency variable
\(\tau\). The equality \(s=\tfrac12+i\tau\) is the Tate half-weight
normalisation, so
\(\tfrac12\mathrm{id}+i\widehat D\) becomes multiplication by \(s\).
\end{proof}

\begin{definition}[archimedean Arakelov disk]
\label{v4-arith:w9-r2:def:Sigma-infty}
\ClaimStatusDefinitional. The \emph{archimedean Arakelov disk} is the
2-dimensional boundary component of \(\overline{\mathrm{Spec}\,\mathbb{Z}}\)
at the archimedean place:
\begin{equation}
\label{v4-arith:w9-r2:eq:Sigma-infty}
\Sigma_{\infty}\;:=\;\{\infty\}\times D^{2}\;\subset\;\overline{\mathrm{Spec}\,\mathbb{Z}},
\end{equation}
where \(D^{2}=\{z\in\mathbb{C}:|z|\leq 1\}\) is the closed unit disk,
equipped with the Faltings 1984 archimedean K\"ahler form
\begin{equation}
\label{v4-arith:w9-r2:eq:omega-arakelov}
\omega_{\infty}\;:=\;\frac{i}{2\pi}\,\frac{dz\wedge d\bar{z}}{(1+|z|^{2})^{2}}
\end{equation}
restricted to \(D^{2}\). The disk carries the
standard U(1)-framing from its complex structure.
\end{definition}

\begin{remark}[role of \(\Sigma_{\infty}\)]
\label{v4-arith:w9-r2:rem:sigma-role}
In the Morishita dictionary (Morishita 2012 \emph{Knots and Primes},
arXiv:0904.3399, \S2) finite primes appear as knots
\(K_{p}\subset\overline{\mathrm{Spec}\,\mathbb{Z}}\). The archimedean
place is structurally a closed surface component rather than a knot:
compactification at \(\infty\) adds a 2-cell. The Wilson observable
attached to \(\infty\) is therefore a Wilson-surface, not a Wilson-line.
The 3D CS bulk couples to the surface via its area form
\eqref{v4-arith:w9-r2:eq:omega-arakelov}; the choice of Faltings
K\"ahler form encodes the Green-metric data of the Arakelov
compactification (Faltings 1984 \emph{Ann.~Math.}~119, 387--424).
\end{remark}

\section{Wilson-Surface Observables in 3D Chern-Simons}
\label{v4-arith:w9-r2:wilson-cs}

The Wilson-surface observable in 3D gauge theory is the higher
categorical analogue of the Wilson-line observable. For a 2D surface
\(\Sigma\subset M^{3}\) in a 3-manifold carrying a gauge connection
\(A\), a Wilson-surface is specified by a 2-representation of the gauge
2-group and evaluates the exponentiated surface holonomy of \(A\)
(Baez--Schreiber 2005 arXiv:hep-th/0412325, \S4; Schreiber--Sati 2011
arXiv:1011.4735, \S2.6). For abelian \(U(1)\) gauge theory, a
Wilson-surface is the exponentiated integral of the curvature 2-form
\(F_{A}=dA\) over \(\Sigma\), weighted by a real charge \(q\):
\begin{equation}
\label{v4-arith:w9-r2:eq:wilson-U1}
W_{q}(\Sigma,A)\;:=\;\exp\!\left(iq\int_{\Sigma}F_{A}\right).
\end{equation}

\subsection{3D Chern-Simons and Wilson-Surface Observables (Witten 1989)}
\label{v4-arith:w9-r2:witten-framework}

Witten 1989 \emph{Commun.~Math.~Phys.}~121, 351--399 (\emph{Quantum
field theory and the Jones polynomial}, originally announced 1988)
established 3D Chern-Simons as a topological gauge theory whose
boundary WZW admits boundary observables specified by knots (Wilson
lines) and by surfaces (Wilson surfaces for higher-form symmetries).
The abelian level-\(k\) action is
\begin{equation}
\label{v4-arith:w9-r2:eq:CS-action}
S_{\mathrm{CS},k}[A]\;=\;\frac{k}{4\pi}\int_{M^{3}}A\wedge dA,
\end{equation}
and the path-integral evaluates correlators of Wilson observables.
Three structural properties of Wilson-surfaces in this framework
are used below.

\begin{proposition}[structural properties of Wilson-surfaces]
\label{v4-arith:w9-r2:prop:wilson-surface}
\ClaimStatusProvedElsewhere. In 3D abelian Chern-Simons at level \(k\)
on a 3-manifold \(M^{3}\):
\begin{enumerate}
\item (Framing dependence.) The Wilson-surface observable
\eqref{v4-arith:w9-r2:eq:wilson-U1} is invariant under smooth
isotopies of \(\Sigma\) preserving the framing.
\item (Charge quantization.) Gauge invariance
requires \(qk\in\mathbb{Z}\) when \(\Sigma\) is closed and \(M^{3}\) is
closed.
\item (Scaling interpretation.) When \(\Sigma=\{p\}\times D^{2}\) sits
over a codimension-1 point \(p\in M^{2}\times\{pt\}\subset M^{3}\) and
\(A\) couples as an area-measurement field, the observable
\eqref{v4-arith:w9-r2:eq:wilson-U1} generates a 1-parameter scaling of
the fibre at \(p\).
\end{enumerate}
\end{proposition}

\begin{proof}
Properties (1) and (2) are Witten 1989, \S3--\S4. Property (3) is the
dual statement: the 2-form \(F_{A}|_{\Sigma}\) integrates to an area
measurement on \(\Sigma\); exponentiation generates a 1-parameter group
of rescalings of that area, realising the scaling flow on the fibre
over \(p\). See Cattaneo--Mnev--Reshetikhin 2014 arXiv:1201.0290, \S5,
for the dual scaling interpretation.
\end{proof}

\subsection{The Archimedean Wilson-Surface}
\label{v4-arith:w9-r2:arch-wilson}

We apply (\Cref{v4-arith:w9-r2:prop:wilson-surface}) with
\(\Sigma=\Sigma_{\infty}\) of \Cref{v4-arith:w9-r2:def:Sigma-infty},
\(M^{3}=\overline{\mathrm{Spec}\,\mathbb{Z}}\), and
\(A=A_{k^{\mathrm{Ar}}}\) the level-\(k^{\mathrm{Ar}}\) arithmetic
Chern-Simons connection of \Cref{v4-arith:3d-acs:conj:qu-acs}. The
area-measurement field on \(\Sigma_{\infty}\) is the pullback of
\eqref{v4-arith:w9-r2:eq:omega-arakelov}.

\begin{definition}[area 2-representation]
\label{v4-arith:w9-r2:def:sigma-area}
\ClaimStatusDefinitional. The \emph{area-scaling 2-representation}
\(\sigma_{\mathrm{area}}\) is the 1-dimensional abelian 2-representation
of the \(U(1)\) gauge 2-group defined by the character
\(\theta\mapsto\exp(i\theta\cdot\mathrm{vol}(\Sigma_{\infty},\omega_{\infty}))\),
where \(\mathrm{vol}(\Sigma_{\infty},\omega_{\infty})=\int_{\Sigma_{\infty}}\omega_{\infty}\)
is the Faltings-normalised Arakelov volume of \(\Sigma_{\infty}\).
\end{definition}

\begin{definition}[archimedean Wilson-surface \(F_{\infty}^{\mathrm{WS}}\)]
\label{v4-arith:w9-r2:def:F-infty-WS}
\ClaimStatusDefinitional. Put
\begin{equation}
\label{v4-arith:w9-r2:eq:F-infty-WS}
F_{\infty}^{\mathrm{WS}}
\;:=\;W_{\sigma_{\mathrm{area}}}(\Sigma_{\infty},A_{k^{\mathrm{Ar}}})
\;=\;\exp\!\left(i\int_{\Sigma_{\infty}}
\omega_{\infty}\cdot F_{A_{k^{\mathrm{Ar}}}}\right),
\end{equation}
the bulk Wilson-surface observable of the 3D quantum arithmetic
Chern-Simons on \(\overline{\mathrm{Spec}\,\mathbb{Z}}\) with label
\(\sigma_{\mathrm{area}}\).
\end{definition}

\begin{remark}[conditional domain]
\label{v4-arith:w9-r2:rem:cond-wilson}
The existence of the bulk 3D quantum arithmetic Chern-Simons measure
at level \(k^{\mathrm{Ar}}\) is the conjectural content of
\Cref{v4-arith:3d-acs:conj:qu-acs}. \(F_{\infty}^{\mathrm{WS}}\) is defined here as an observable of that
(conjectural) measure. The construction is unconditional as a
structural formula; the well-definedness of the measure against
which expectation values are taken is conditional.
\end{remark}

\subsection{Source Comparison for the Wilson-Surface Construction}
\label{v4-arith:w9-r2:ah:wilson}

The three source classes are disjoint: Witten 1989
(\emph{Comm.~Math.~Phys.}~121) for Wilson observables in 3D
Chern-Simons; Baez--Schreiber 2005 for the 2-representation language;
Faltings 1984 for the Arakelov K\"ahler form.

Sign and convention: \(A\) is a real connection, \(F_{A}=dA\) is a real
2-form, and the \(i\)-exponential gives a unitary observable. The
Faltings K\"ahler form is positive on \(\Sigma_{\infty}\) so the area
integral is positive real, consistent with
\Cref{v4-arith:w9-r2:lem:mellin}.

The observable category is unitary operators on the boundary Hilbert
space of \(\partial\mathrm{ACS}^{\mathrm{qu}}\) at \(\Sigma_{\infty}\),
which is the archimedean Fock space
\(\mathcal{H}^{(\infty)}=\mathcal{V}^{\mathrm{prim}}_{\mathrm{Heis},\infty}\)
of Ch.~\ref{v4-arith:w5-b:hp-operator}; \(F_{\infty}^{\mathrm{WS}}\)
lands in \(\mathrm{End}(\mathcal{H}^{(\infty)})\).

The bulk measure of \Cref{v4-arith:3d-acs:conj:qu-acs} is conditional;
the observable as a structural formula is unconditional. The residual
is bulk measure existence.

\(F_{\infty}^{\mathrm{WS}}\) is functorial in the gauge 2-group along
the \(U(1)\)-inclusion. Non-abelian gauge groups would add
Stokes-boundary corrections (Baez--Schreiber 2005, \S5); the abelian
case bypasses these. The identification of \(F_{\infty}^{\mathrm{WS}}\)
with the Arakelov scaling derivative is the content of
\Cref{v4-arith:w9-r2:thm:WS-equals-scaling}.

\section{Construction of \(F_{\infty}\)}
\label{v4-arith:w9-r2:construction-F-infty}

We construct \(F_{\infty}\) via Arakelov scaling on the archimedean
fibre, independent of the conditional Wilson-surface realisation, and
prove that the two constructions agree.

\begin{definition}[archimedean Frobenius as scaling exponential]
\label{v4-arith:w9-r2:def:F-infty}
\ClaimStatusDefinitional. The \emph{archimedean Frobenius} is the
operator
\begin{equation}
\label{v4-arith:w9-r2:eq:F-infty-def}
F_{\infty}\;:=\;\exp\bigl(-\Theta_{\infty}\bigr)
\;=\;\exp\!\left(-\tfrac12\mathrm{id}-i\widehat D\right)
\end{equation}
acting on the archimedean Fock factor
\(\mathcal{H}^{(\infty)}\subset\mathcal{V}^{\mathrm{prim}}_{\mathrm{Heis}}\),
realised on \(L^{2}(\mathbb{R},dt)\) via the Weil--Petersson isometry
\(\Upsilon_{\infty}\) of
Ch.~\ref{v4-arith:w5-a:chapter}.
Equivalently, on the Mellin side it is multiplication by \(e^{-s}\)
on the critical line by \Cref{v4-arith:w9-r2:lem:mellin}.
\end{definition}

\begin{remark}[heat-flow form]
\label{v4-arith:w9-r2:rem:heat}
Alternative presentation: \(F_{\infty}=e^{-\Theta_{\infty}}\), where
\(\Theta_{\infty}=\tfrac12\mathrm{id}+i\widehat D\) and
\(\widehat D\) is the Connes 1999 multiplicative scaling generator of
\Cref{v4-arith:w9-r2:def:scaling}. Under the Mellin unitarity of
\Cref{v4-arith:w9-r2:lem:mellin}, \(F_{\infty}\) becomes the
Mellin-multiplier operator with symbol \(e^{-s}\).
This is the Frobenius eigenvalue convention. It must be kept distinct
from the external variable in a multiplicative Euler-character
determinant \(\det_{\infty}(1-e^{-s}F\mid H^{\bullet}_{\mathrm{ch}})\).
\end{remark}

\begin{theorem}[Wilson-surface \(=\) Arakelov scaling]
\label{v4-arith:w9-r2:thm:WS-equals-scaling}
\ClaimStatusProvedHereConditional \emph{(on
\Cref{v4-arith:3d-acs:conj:qu-acs})}. Under the conjectural bulk
measure of quantum arithmetic Chern-Simons, the Wilson-surface
observable \(F_{\infty}^{\mathrm{WS}}\) of
\Cref{v4-arith:w9-r2:def:F-infty-WS} coincides with \(F_{\infty}\) of
\Cref{v4-arith:w9-r2:def:F-infty} on the archimedean Fock factor:
\begin{equation}
\label{v4-arith:w9-r2:eq:WS-equals-scaling}
F_{\infty}^{\mathrm{WS}}\;=\;F_{\infty}\qquad\text{in}\ \mathrm{End}(\mathcal{H}^{(\infty)}).
\end{equation}
\end{theorem}

\begin{proof}
Under \Cref{v4-arith:3d-acs:conj:qu-acs}, the boundary state Hilbert
space at \(\Sigma_{\infty}\) is
\(\mathcal{H}^{(\infty)}\cong\mathcal{V}^{\mathrm{prim}}_{\mathrm{Heis},\infty}\);
this is the boundary WZW Hilbert space axiom (1) of the conjecture.
On this Hilbert space the area-coupled Chern-Simons connection
\(A_{k^{\mathrm{Ar}}}\) acts by the level-\(k^{\mathrm{Ar}}\) primary
current; the curvature \(F_{A}=dA\) acts by the derivative of the
primary current in the normal direction. Integration over
\(\Sigma_{\infty}\) against the Arakelov K\"ahler form
\eqref{v4-arith:w9-r2:eq:omega-arakelov} is precisely the archimedean
scaling-generator computation: the pullback of \(\omega_{\infty}\) is
proportional to \(d^{\times}|x|\wedge d\arg(x)\) on the unit-norm locus,
and integrating against the normal curvature gives the normal
operator \(\Theta_{\infty}\), built from the generator
\(\widehat{D}\) of \Cref{v4-arith:w9-r2:def:scaling}, up to the
normalization \(1/(2\pi)\cdot\mathrm{vol}(\Sigma_{\infty},\omega_{\infty})=1\)
fixed by Faltings 1984 (\S2--\S3). Exponentiation gives
\eqref{v4-arith:w9-r2:eq:WS-equals-scaling}. The residual is the
well-definedness of the bulk measure, which is exactly
\Cref{v4-arith:3d-acs:conj:qu-acs}.
\end{proof}

\begin{remark}[geometric content of the identification]
\label{v4-arith:w9-r2:rem:geom-content}
The identification is structurally the same mechanism by which boundary
\(L_{0}\) in 2D CFT equals the bulk radial Hamiltonian in the 3D
Chern-Simons dual: radial quantisation of the 3D bulk turns
\(\partial_{t}\) (where \(t\) is the radial coordinate) into
\(L_{0}^{\mathrm{Vir}}\) on the boundary Hilbert space (Witten 1989,
\S4; Elitzur--Moore--Schwimmer--Seiberg 1989 \emph{Nucl.~Phys.}~B326,
108--134). On the arithmetic side, the radial coordinate is the
Arakelov log-valuation parameter \(t=\log|x|_{\infty}\), and the
boundary Hamiltonian is the self-adjoint generator \(\widehat{D}\) of
\Cref{v4-arith:w9-r2:def:scaling}.
\end{remark}

\subsection{Convention and Source Comparison for the Construction}
\label{v4-arith:w9-r2:ah:construction}

Three source classes are disjoint: Arakelov scaling (Arakelov 1974,
Faltings 1984, Connes 1999); bulk Wilson-surface (Witten 1989,
Baez--Schreiber 2005); radial quantisation identification
(Elitzur--Moore--Schwimmer--Seiberg 1989).

Sign conventions: \(\widehat D=-i\partial_t\) is self-adjoint on
\(L^{2}(\mathbb R,dt)\), \(\Theta_{\infty}=\tfrac12\mathrm{id}+i\widehat D\),
and \(F_{\infty}=e^{-\Theta_{\infty}}\). Thus \(F_{\infty}\) is
\(e^{-1/2}\) times the unitary scaling flow. This is the Deninger
normalisation: a zero on the critical line,
\(\rho=\tfrac12+i\tau\), gives
\(|e^{-\rho}|=e^{-1/2}\).

The ambient Hilbert space is
\(L^{2}(\mathbb{R},dt)\simeq L^{2}(\mathrm{Re}\,s=\frac{1}{2},|ds|/(2\pi))\),
unitarily identified with \(\mathcal{H}^{(\infty)}\) via the
Weil--Petersson isometry of Ch.~\ref{v4-arith:w5-a:chapter}. The
unitary part of the observable is the exponential of the self-adjoint
generator \(\widehat D\); the Frobenius observable includes the Tate
factor \(e^{-1/2}\).

The quantum ACS bulk measure is required for Wilson-surface evaluation;
this is the residual of \Cref{v4-arith:3d-acs:conj:qu-acs}.

Under the Tate-pole cancellation \(\tfrac{1}{2}s(s-1)\) of
Ch.~\ref{v4-arith:w8-a:H-structure}, the \(F_{\infty}\)-eigenvalues at
\(s=0,1\) cancel and the critical-line eigenvalues survive, consistent
with \Cref{v4-arith:w8-a:prop:H1-structure}.

For the disk form in \eqref{v4-arith:w9-r2:eq:omega-arakelov},
\(\int_{D^{2}}\omega_{\infty}=1/2\). A unit-area convention would replace
\(\omega_{\infty}\) by \(2\omega_{\infty}\), or else use the full
projective archimedean fibre. The Wilson-surface comparison below uses
the scaling class, so no numerical area normalisation is inserted.

\section{Conjugacy to Virasoro \(L_{0}^{\mathrm{Ar}}\)}
\label{v4-arith:w9-r2:conjugacy}

The third description of \(F_{\infty}\) is as the Mellin-conjugate of
the arithmetic Virasoro zero-mode \(L_{0}^{\mathrm{Ar}}\). Chapter~\ref{v4-arith:rh-reformulation}
identified \(L_{0}^{\mathrm{Ar}}\) on the primary subspace
\(\mathcal{V}^{\mathrm{prim}}_{\mathrm{Heis},\mathrm{prim}}\) as the
Fock/Virasoro conformal-energy operator of the formal arithmetic Fock.
Chapter~\ref{v4-arith:w8-a:machinery:ch23} recorded that
\(L_{0}^{\mathrm{Ar}}\) is distinct from \(\widehat{D}\),
\(H_{\mathrm{HP}}\), and \(\Theta_{\xi}\). The
Weil--Petersson isometry \(\Upsilon_{\infty}\) identifies the
normal archimedean conformal-weight operator
\(\Theta_{\infty}=\tfrac12\mathrm{id}+i\widehat D\) with
\(L_{0}^{\mathrm{Ar}}\) on the archimedean primary fibre, while
\(H_{\mathrm{HP}}\) and \(\Theta_{\xi}\) remain distinct global
operators on the full \(\mathcal{V}^{\mathrm{prim}}\).

\subsection{The Weil--Petersson isometry \(\Upsilon_{\infty}\)}
\label{v4-arith:w9-r2:WP-isometry}

Chapter~\ref{v4-arith:w5-a:chapter} constructed a Weil--Petersson
unitary isometry
\begin{equation}
\label{v4-arith:w9-r2:eq:WP-isometry}
\Upsilon_{\infty}\;\colon\;
\mathcal{V}^{\mathrm{prim}}_{\mathrm{Heis},\mathrm{prim},\infty}
\xrightarrow{\ \sim\ }L^{2}(\mathbb{R}^{\times}_{+},d^{\times}x)
\end{equation}
on the archimedean fibre restricted to the primary sector. Under
\(\Upsilon_{\infty}\):
\begin{itemize}
\item The Petersson inner product on the primary Fock becomes the
standard \(L^{2}\) inner product.
\item The Arakelov scaling generator \(\widehat{D}\) of
\Cref{v4-arith:w9-r2:def:scaling} becomes multiplication by
\(\tau\) on the Mellin transform.
\item The arithmetic conformal-energy operator \(L_{0}^{\mathrm{Ar}}\)
is required, in the determinant-trace hypothesis below, to become
the Mellin multiplier \(s=\tfrac12+i\tau\) on the primary sector.
\end{itemize}

The third property is the essential content of
\Cref{v4-arith:w9-r2:lem:L0-Mellin}.

\begin{hypothesis}[\(\Upsilon_{\infty}\) realises \(L_{0}^{\mathrm{Ar}}\) as Mellin multiplication]
\label{v4-arith:w9-r2:lem:L0-Mellin}
\ClaimStatusConjectured. On the primary subspace
\(\mathcal{V}^{\mathrm{prim}}_{\mathrm{Heis},\mathrm{prim},\infty}\),
\begin{equation}
\label{v4-arith:w9-r2:eq:L0-Mellin}
\Upsilon_{\infty}\,L_{0}^{\mathrm{Ar}}\,\Upsilon_{\infty}^{-1}
\;=\;\text{multiplication by}\ s
\quad\text{on}\ L^{2}(\mathrm{Re}\,s=1/2).
\end{equation}
\end{hypothesis}

\begin{remark}[spectral type obstruction]
The ordinary archimedean Fock \(L_{0}\) has discrete spectrum
\(\mathbb{Z}_{\geq0}\), whereas multiplication by \(s\) on the Mellin
line has continuous spectrum. Therefore
\Cref{v4-arith:w9-r2:lem:L0-Mellin} is an additional
determinant-trace assertion. It is not a consequence of the primary
bar Mellin character \(\sum_{m\geq1}m^{-s}=\zeta(s)\), and it is not
supplied by the determinant residue \(r_{\zeta}=2\).
\end{remark}

\subsection{The conjugacy theorem}
\label{v4-arith:w9-r2:conjugacy-theorem}

\begin{theorem}[archimedean Frobenius \(=\) Virasoro \(L_{0}\) exponential]
\label{v4-arith:w9-r2:thm:conjugacy}
\ClaimStatusProvedHereConditional \emph{(on
\Cref{v4-arith:3d-acs:conj:qu-acs},
\Cref{v4-arith:rh:conj:primary-Virasoro}, and
\Cref{v4-arith:w9-r2:lem:L0-Mellin})}. On the archimedean primary
Fock \(\mathcal{V}^{\mathrm{prim}}_{\mathrm{Heis},\mathrm{prim},\infty}\),
\begin{equation}
\label{v4-arith:w9-r2:eq:conjugacy}
F_{\infty}\;=\;\Upsilon_{\infty}^{-1}\,e^{-L_{0}^{\mathrm{Ar}}}\,\Upsilon_{\infty}.
\end{equation}
Equivalently, under \(\Upsilon_{\infty}\), the archimedean Frobenius is
the Mellin multiplier \(e^{-s}\), and the Virasoro zero-mode is
multiplication by \(s\). Together with
\Cref{v4-arith:w9-r2:thm:WS-equals-scaling}, the three constructions
\(F_{\infty}\), \(F_{\infty}^{\mathrm{WS}}\), and
\(\Upsilon_{\infty}^{-1}\,e^{-L_{0}^{\mathrm{Ar}}}\,\Upsilon_{\infty}\)
agree.
\end{theorem}

\begin{proof}
By \Cref{v4-arith:w9-r2:lem:mellin},
\(\Upsilon_{\infty}\,\Theta_{\infty}\,\Upsilon_{\infty}^{-1}=s\).
By \Cref{v4-arith:w9-r2:lem:L0-Mellin},
\(\Upsilon_{\infty}\,L_{0}^{\mathrm{Ar}}\,\Upsilon_{\infty}^{-1}=s\).
The two Mellin-images coincide; therefore
\(\Theta_{\infty}=L_{0}^{\mathrm{Ar}}\) after \(\Upsilon_{\infty}\)-conjugation.
Exponentiating gives \(F_{\infty}=e^{-\Theta_{\infty}}\) conjugate to
\(e^{-L_{0}^{\mathrm{Ar}}}\) via \(\Upsilon_{\infty}\).
\end{proof}

\begin{remark}[three operators, one Mellin symbol]
\label{v4-arith:w9-r2:rem:one-symbol}
The operator dictionary of Chapter~\ref{v4-arith:rh:rem:spectral-convention}
distinguishes four operators: \(L_{0}^{\mathrm{Vir}}\), \(N_{\mathrm{BC}}\),
\(H_{\mathrm{HP}}\), \(\Theta_{\xi}\). The present theorem identifies
\(\Theta_{\infty}\), not \(\widehat D\), with \(L_{0}^{\mathrm{Ar}}\)
on the \emph{archimedean} primary fibre only, under Weil--Petersson
conjugation. This is a local archimedean statement; it does not
identify \(L_{0}^{\mathrm{Vir}}\) with \(H_{\mathrm{HP}}\) on the full
\(\mathcal{V}^{\mathrm{prim}}\), where the finite primes contribute a
separate \(N_{\mathrm{BC}}\)-action and the Hilbert--P\'olya operator
\(H_{\mathrm{HP}}\) is the conjectural determinant-trace operator attached to
the assembled zero distribution. The Polyakov divergence of
\Cref{v4-arith:rh:rem:Virasoro-reconciliation} is at the full global
level, not at the archimedean-fibre level of
\Cref{v4-arith:w9-r2:thm:conjugacy}.
\end{remark}

\subsection{Convention and Source Comparison for the Conjugacy}
\label{v4-arith:w9-r2:ah:conjugacy}

Source classes are disjoint: Arakelov scaling (\(\widehat{D}\)) from
Arakelov 1974 and Connes 1999; Fock-Virasoro (\(L_{0}^{\mathrm{Ar}}\))
from Segal 1981 and Kac--Raina 1987.

The self-adjoint generator \(\widehat D\) has Mellin-image \(\tau\).
The conformal-weight operator
\(\Theta_{\infty}=\tfrac12\mathrm{id}+i\widehat D\) and
\(L_{0}^{\mathrm{Ar}}\) have Mellin-image \(s=\tfrac12+i\tau\).

The ambient category is the Hilbert space
\(L^{2}(\mathrm{Re}\,s=\tfrac{1}{2},|ds|/(2\pi))\). The unitary
operator is \(e^{-i\widehat D}\). The Frobenius is
\(e^{-\Theta_{\infty}}\).

Primary-Virasoro reconciliation (\Cref{v4-arith:rh:conj:primary-Virasoro})
is required for \(L_{0}^{\mathrm{Ar}}\) to be an honest centrally-extended
Virasoro on the primary core. This is the residual.

Under Arakelov symmetry \(\lambda\mapsto\lambda^{-1}\), \(F_{\infty}\)
inverts; the Mellin-side involution \(s\mapsto 1-s\) also inverts, compatible
with Koszul self-duality of \Cref{v4-arith:thm:P1-resolution}.

The equality \(\Theta_{\infty}=L_{0}^{\mathrm{Ar}}\) under
\(\Upsilon_{\infty}\) is exactly the Mellin realisation hypothesis
\Cref{v4-arith:w9-r2:lem:L0-Mellin}, together with
\Cref{v4-arith:w9-r2:lem:mellin}.

\section{Spectral Identification with \(\xi_{\mathbb{R}}\)-zeros}
\label{v4-arith:w9-r2:spectral}

The conjugacy of
\Cref{v4-arith:w9-r2:thm:conjugacy} transfers the spectrum of
\(L_{0}^{\mathrm{Ar}}\) on the primary Fock to the spectrum of
\(F_{\infty}\) on the archimedean Fock factor. Under the conjugate
Hilbert--P\'olya hypothesis H-P\(^{\mathrm{Ar}}\) of
\Cref{v4-arith:rh:hyp:HPAr}, the middle-degree descent
\(H^{1}_{\mathrm{ch}}\) receives a regularised \(F\)-action whose
determinant divisor is the exponential image of the non-trivial-zero
divisor of \(\xi_{\mathbb{R}}\).

\subsection{\(L_{0}^{\mathrm{Ar}}\) on the primary Fock and the \(\xi\)-image}
\label{v4-arith:w9-r2:L0-spectrum}

Chapter~\ref{v4-arith:rh-reformulation} identified \(\xi_{\mathbb{R}}\)
as the character of \(\mathcal{V}^{\mathrm{prim}}\) via Tate 1950
completion of the Heisenberg bar-character \(\zeta\). The character
encodes the Tate-shifted Mellin determinant after Petersson
normalisation. Under the H-P\(^{\mathrm{Ar}}\) realisation, the normal
operator \(\Theta_{\xi}\) has regularised determinant divisor equal to
\(\mathrm{div}\,\xi_{\mathbb R}\); after zero-placement this is the
ordinary ordinate spectrum of its self-adjoint part.

\begin{lemma}[\(L_{0}^{\mathrm{Ar}}\) determinant divisor via HP zeros]
\label{v4-arith:w9-r2:lem:L0-zeros}
\ClaimStatusProvedHereConditional \emph{(on
\Cref{v4-arith:rh:hyp:HPAr} and
\Cref{v4-arith:w9-r2:lem:L0-Mellin})}. Under these hypotheses, the
regularised determinant divisor of the Mellin-image
\(s=\Upsilon_{\infty}\,L_{0}^{\mathrm{Ar}}\,\Upsilon_{\infty}^{-1}\)
on the middle-degree descent \(H^{1}_{\mathrm{ch}}\) is
\begin{equation}
\label{v4-arith:w9-r2:eq:L0-eigenvalues}
\operatorname{div}_{\infty}\bigl(\Upsilon_{\infty}\,L_{0}^{\mathrm{Ar}}\,
\Upsilon_{\infty}^{-1}\mid H^{1}_{\mathrm{ch}}\bigr)
\;=\;\operatorname{div}\xi_{\mathbb{R}}.
\end{equation}
\end{lemma}

\begin{proof}
Under H-P\(^{\mathrm{Ar}}\), \(\Theta_{\xi}=\tfrac{1}{2}\mathrm{id}+iH_{\mathrm{HP}}\)
regularised-determinant identity
\(\det_{\infty}(s\cdot\mathrm{id}-\Theta_{\xi})=\xi_{\mathbb{R}}(s)\) of
\Cref{v4-arith:rh:eq:HP-determinant} says precisely that
\(\Theta_{\xi}\) has determinant divisor \(\operatorname{div}\xi_{\mathbb R}\).
The Mellin multiplier-image of \(L_{0}^{\mathrm{Ar}}\) by
\Cref{v4-arith:w9-r2:lem:L0-Mellin} is the same Tate-shifted
Mellin variable \(s\).
\end{proof}

\subsection{\(F_{\infty}\) determinant as the exponential zero image}
\label{v4-arith:w9-r2:F-eigenvalues}

\begin{theorem}[determinant-divisor identification for \(F_{\infty}\)]
\label{v4-arith:w9-r2:thm:spectral-identification}
\ClaimStatusProvedHereConditional \emph{(on
\Cref{v4-arith:rh:hyp:HPAr},
\Cref{v4-arith:rh:conj:primary-Virasoro},
\Cref{v4-arith:3d-acs:conj:qu-acs}, and
\Cref{v4-arith:w9-r2:lem:L0-Mellin})}. The regularised determinant
divisor of
\(F_{\infty}\) on the middle-degree archimedean descent of the chiral
cohomology is the pushforward of \(\operatorname{div}\xi_{\mathbb R}\)
under \(\rho\mapsto e^{-\rho}\):
\begin{equation}
\label{v4-arith:w9-r2:eq:F-eigenvalues}
\operatorname{div}_{\infty}\bigl(F_{\infty}\mid H^{1}_{\mathrm{ch}}\bigr)
\;=\;(e^{-\bullet})_{*}\operatorname{div}\xi_{\mathbb{R}},
\end{equation}
with multiplicities summed over fibres of the exponential map. After
zero-placement this divisor is represented by the spectral measure of
the normal operator \(F_{\infty}\).
\end{theorem}

\begin{proof}
By \Cref{v4-arith:w9-r2:thm:conjugacy},
\(F_{\infty}=\Upsilon_{\infty}^{-1}e^{-L_{0}^{\mathrm{Ar}}}\Upsilon_{\infty}\)
on the primary archimedean Fock. By \Cref{v4-arith:w9-r2:lem:L0-zeros},
the determinant divisor of the Mellin image of \(L_{0}^{\mathrm{Ar}}\)
on \(H^{1}_{\mathrm{ch}}\) is \(\operatorname{div}\xi_{\mathbb R}\).
Holomorphic functional calculus pushes a determinant divisor forward
under \(z\mapsto e^{-z}\), giving
\eqref{v4-arith:w9-r2:eq:F-eigenvalues}.
\end{proof}

\begin{proposition}[multiplicative character does not replace Hadamard]
\label{v4-arith:w9-r2:cor:character-factorisation}
\ClaimStatusProvedHere. The multiplicative determinant
\begin{equation}
\label{v4-arith:w9-r2:eq:character-factorisation}
\det\nolimits_{\infty}\bigl(1-e^{-s}F_{\infty}\mid H^{1}_{\mathrm{ch}}\bigr)
\;=\;\prod_{\rho}\bigl(1-e^{-s}e^{-\rho}\bigr)
\end{equation}
is not the Hadamard product for \(\xi_{\mathbb R}(s)\). It vanishes at
the logarithmic congruence classes \(s=-\rho+2\pi i\mathbb Z\), not at
\(s=\rho\). The zeta determinant is the additive determinant
\[
\det\nolimits_{\infty}(s\mathrm{id}-\Theta_{\xi}),
\]
not the multiplicative Euler-character determinant
\eqref{v4-arith:w9-r2:eq:character-factorisation}.
\end{proposition}

\begin{proof}
The zeros of the right hand side of
\eqref{v4-arith:w9-r2:eq:character-factorisation} satisfy
\(e^{-s}e^{-\rho}=1\), hence \(s=-\rho+2\pi i k\). Hadamard's product
has local factors \((1-s/\rho)e^{s/\rho}\), with zeros at
\(s=\rho\). Exponentiating the spectrum therefore changes the divisor.
Only the additive regularised determinant keeps the zero divisor.
\end{proof}

\subsection{Convention and Source Comparison for the Spectral Identification}
\label{v4-arith:w9-r2:ah:spectral}

The Hilbert--P\'olya conjecture (Odlyzko 1987 \emph{Math.~Comp.}~48;
Berry--Keating 1999 \emph{SIAM Rev.}~41) and Arakelov scaling are
independent.

The \(F_{\infty}\)-eigenvalue attached to \(\rho\) is \(e^{-\rho}\).
This is the multiplicative Frobenius notation. It does not replace
the additive determinant \(\det_{\infty}(s\mathrm{id}-\Theta_{\xi})\).

The ambient category is a Hilbert space with regularised trace and
regularised determinant on nuclear operators (Quillen 1985
\emph{Funct.~Anal.~Appl.}~19).

The hypothesis H-P\(^{\mathrm{Ar}}\) is required for the
Tate-shifted \(L_{0}^{\mathrm{Ar}}\) determinant on
\(H^{1}_{\mathrm{ch}}\) to have divisor \(\mathrm{div}\,\xi_{\mathbb R}\).

The zero divisor is closed under \(\rho\mapsto 1-\rho\) by the
functional equation. The \(F_{\infty}\) determinant divisor is its
exponential image and carries the induced involution.

The \(F_{\infty}\) Mellin multiplier is \(e^{-s}\). The additive
operator whose determinant has divisor \(\mathrm{div}(\xi)\) is
\(\Theta_{\xi}\), not \(F_{\infty}\).

\section{Archimedean Purity is Equivalent to A-TP}
\label{v4-arith:w9-r2:purity-ATP}

Deninger axiom (6) of \Cref{v4-arith:w8-a:deninger-wishlist} requires
that the Frobenius acts on \(H^{1}_{\Delta}\) with weight \(1/2\):
\(|F|\)-eigenvalues on \(H^{1}\) have absolute value \(e^{-1/2}\) for
the normalisation \(F=\exp(-\Theta_{\xi})\) with
\(\mathrm{Re}(\Theta_{\xi})=1/2\). Equivalently, the unitarised
operator \(e^{1/2}F\) has spectral radius one. The archimedean
component of this purity statement is the claim
\(|F_{\infty}|\)-weight \(=1/2\). This is equivalent to archimedean
Tate positivity A-TP by
\Cref{v4-arith:w9-r2:thm:purity-ATP-RH}.

\subsection{The \(|F_{\infty}|\)-weight \(=1/2\) statement}
\label{v4-arith:w9-r2:F-weight-statement}

\begin{definition}[archimedean \(F_{\infty}\)-weight]
\label{v4-arith:w9-r2:def:F-weight}
\ClaimStatusDefinitional. The \emph{archimedean \(F_{\infty}\)-weight
on \(H^{1}_{\mathrm{ch}}\)} is the common value \(w\in\mathbb{R}\) such
that every regularised eigenvalue \(\mu\) of \(F_{\infty}\) on
\(H^{1}_{\mathrm{ch}}\) satisfies \(|\mu|=e^{-w}\), provided such a
common value exists; otherwise the \(F_{\infty}\)-weight is not well-defined.
\end{definition}

\begin{remark}[well-definedness]
\label{v4-arith:w9-r2:rem:well-def}
By \Cref{v4-arith:w9-r2:thm:spectral-identification}, the
\(F_{\infty}\) eigenvalues are
\(\{e^{-\rho}:\xi_{\mathbb{R}}(\rho)=0\}\). The modulus
\(|e^{-\rho}|=e^{-\mathrm{Re}(\rho)}\) equals \(e^{-1/2}\) for all zeros
exactly when all zeros satisfy \(\mathrm{Re}(\rho)=1/2\), i.e.\ exactly
when RH holds. Thus the \(F_{\infty}\)-weight is well-defined and
equals \(1/2\) iff RH.
\end{remark}

\subsection{A-TP archimedean purity equivalence}
\label{v4-arith:w9-r2:ATP-equivalence}

\begin{theorem}[\(F_{\infty}\)-purity \(\Leftrightarrow\) A-TP \(\Leftrightarrow\) RH]
\label{v4-arith:w9-r2:thm:purity-ATP-RH}
\ClaimStatusProvedHereConditional \emph{(on
\Cref{v4-arith:rh:hyp:HPAr},
\Cref{v4-arith:rh:conj:primary-Virasoro},
\Cref{v4-arith:3d-acs:conj:qu-acs},
\Cref{v4-arith:w5-a:thm:ATP-equals-weil-positivity})}. The following
four statements are equivalent:
\begin{enumerate}
\item The \(F_{\infty}\)-weight on \(H^{1}_{\mathrm{ch}}\) is well-defined
and equals \(1/2\): every regularised eigenvalue of \(F_{\infty}\) on
\(H^{1}_{\mathrm{ch}}\) has modulus \(e^{-1/2}\).
\item Deninger archimedean purity axiom (6): the scaling Frobenius
at \(\infty\) acts on \(H^{1}\) at real part \(1/2\).
\item A-TP (\Cref{v4-arith:w5-a:conj:archimedean-tate-positivity}):
archimedean Tate positivity on admissible Paley--Wiener test data.
\item Riemann Hypothesis: all non-trivial zeros \(\rho\) of \(\zeta\)
satisfy \(\mathrm{Re}(\rho)=1/2\).
\end{enumerate}
\end{theorem}

\begin{proof}
(1) \(\Leftrightarrow\) (4): by \Cref{v4-arith:w9-r2:rem:well-def},
\(|e^{-\rho}|=e^{-1/2}\) for all zeros iff \(\mathrm{Re}(\rho)=1/2\) for
all zeros, which is RH.

(2) \(\Leftrightarrow\) (1): Deninger axiom (6) at \(\infty\) requires
the scaling Frobenius to act on \(H^{1}\) at real part \(1/2\);
\(F_{\infty}\) is the scaling Frobenius at \(\infty\) via
\Cref{v4-arith:w9-r2:thm:WS-equals-scaling} and
\Cref{v4-arith:w9-r2:rem:heat}. The real-part \(1/2\) condition for
\(\Theta_{\xi}\)-eigenvalues translates under the exponential
\(F_{\infty}=e^{-\Theta_{\xi}^{\mathrm{arch}}}\) to
\(|F_{\infty}|\)-eigenvalue-modulus \(e^{-1/2}\).

(3) \(\Leftrightarrow\) (4): this is
\Cref{v4-arith:w5-a:thm:ATP-equals-weil-positivity}: A-TP is
equivalent to Weil positivity for the number-field explicit formula,
which is classically equivalent to RH (Weil 1952, Bombieri 2000).

By transitivity, (1), (2), (3), (4) are equivalent.
\end{proof}

\begin{corollary}[three-construction RH equivalence via \(F_{\infty}\)]
\label{v4-arith:w9-r2:cor:three-constructions}
\ClaimStatusProvedHereConditional \emph{(on
\Cref{v4-arith:w9-r2:thm:purity-ATP-RH})}. The archimedean Frobenius
\(F_{\infty}\) provides the archimedean-only formulation of RH that
aligns with the three constructions of \Cref{v4-arith:w8-a:cor:rh-construction}:
\begin{itemize}
\item \emph{Deninger construction at \(\infty\)}: prove \(F_{\infty}\)-weight
\(=1/2\) on \(H^{1}_{\mathrm{ch}}\).
\item \emph{3D-ACS construction at \(\infty\)}: prove the
\(\sigma_{\mathrm{area}}\)-Wilson-surface observable has
weight-\(1/2\) spectrum.
\item \emph{A-TP archimedean construction}: prove archimedean Tate
positivity.
\end{itemize}
All three are equivalent by \Cref{v4-arith:w9-r2:thm:purity-ATP-RH}.
\end{corollary}

\begin{remark}[the archimedean-only portion of RH]
\label{v4-arith:w9-r2:rem:arch-only}
The full global Frobenius \(F\) combines \(F_{\infty}\) with per-prime
\(\psi_{p}\) via restricted tensor (\Cref{v4-arith:w8-a:def:F-global}).
The per-prime \(\psi_{p}\)-eigenvalues on
\(\mathcal{H}^{(p)}_{\mathrm{BC}}\) have modulus \(p^{-r/2}\) on the
weight-\(r\) subspace, already pure by the Bost--Connes Adams structure
(Borger 2009 arXiv:0906.3146, \S8; Connes--Marcolli 2008 AMS Colloq.~55,
\S3). The archimedean portion \(F_{\infty}\) is the only component
where weight-\(1/2\) purity is non-trivial: it reduces to RH by the
above theorem. This explains the concentration of the Riemann
Hypothesis in the archimedean fibre: it is equivalent to archimedean
purity for a Frobenius whose finite-place components are already pure
by construction.
\end{remark}

\subsection{Convention and Source Comparison for the Purity Equivalence}
\label{v4-arith:w9-r2:ah:purity-ATP}

Source classes are disjoint: Deninger 1998 (purity axiom), Weil 1952
(explicit formula and Weil positivity), Ch.~\ref{v4-arith:w5-a:chapter}
(A-TP formulation).

Sign: \(F_{\infty}=e^{-\Theta_{\xi}^{\mathrm{arch}}}\) gives
\(|F_{\infty}|=e^{-\mathrm{Re}(\Theta_{\xi})}\); Deninger purity
requires \(\mathrm{Re}(\Theta_{\xi})=1/2\), consistent.

The ambient category is self-adjoint operators on a Hilbert space with
regularised spectra.

H-P\(^{\mathrm{Ar}}\) is required for the realisation of \(\Theta_{\xi}\).
This is the residual.

Functoriality under \(\rho\leftrightarrow 1-\rho\) (functional equation)
and under Arakelov scaling (archimedean involution) are compatible.

The equivalence A-TP \(\Leftrightarrow\) Weil positivity
\(\Leftrightarrow\) RH is classical; the equivalence
(1)\(\Leftrightarrow\)(2) is direct from the exponential identification
\(F_{\infty}=e^{-\Theta_{\xi}^{\mathrm{arch}}}\).

\section{Residual Obstructions}
\label{v4-arith:w9-r2:residual}

Three residuals remain open from \Cref{v4-arith:w8-a:rem:wilson-infinity}
and \Cref{v4-arith:w8-a:lem:L0-conjugacy}:

\begin{enumerate}
\item \emph{Quantum ACS bulk measure}
(\Cref{v4-arith:3d-acs:conj:qu-acs}). The Wilson-surface observable
\(F_{\infty}^{\mathrm{WS}}\) of
\Cref{v4-arith:w9-r2:def:F-infty-WS} is a well-defined observable of
the conjectural bulk measure; its identification with the
unconditional Arakelov scaling \(F_{\infty}\)
(\Cref{v4-arith:w9-r2:thm:WS-equals-scaling}) is conditional on the
bulk measure. Domain: the profinite renormalisation programme of
Ch.~\ref{v4-arith:3d-acs-E1bar}, a non-formal theorem condition.
\item \emph{Primary-Virasoro reconciliation}
(\Cref{v4-arith:rh:conj:primary-Virasoro}). The conjugacy
\(F_{\infty}=\Upsilon_{\infty}^{-1}e^{-L_{0}^{\mathrm{Ar}}}\Upsilon_{\infty}\)
of \Cref{v4-arith:w9-r2:thm:conjugacy} requires
\(L_{0}^{\mathrm{Ar}}\) to be an honest centrally-extended Virasoro on
the primary core and to satisfy the Mellin realisation hypothesis
\Cref{v4-arith:w9-r2:lem:L0-Mellin}. The latter is a spectral-type
assertion, not a consequence of the Fock grading.
\item \emph{H-P\(^{\mathrm{Ar}}\) realisation}
(\Cref{v4-arith:rh:hyp:HPAr}). The determinant-divisor identity
\[
\operatorname{div}_{\infty}(F_{\infty}\mid H^{1}_{\mathrm{ch}})
=(e^{-\bullet})_{*}\operatorname{div}\xi_{\mathbb R}
\]
of
\Cref{v4-arith:w9-r2:thm:spectral-identification} is conditional on
H-P\(^{\mathrm{Ar}}\) supplying a self-adjoint operator
\(\Theta_{\xi}=\tfrac{1}{2}\mathrm{id}+iH_{\mathrm{HP}}\) with the
non-trivial-zero regularised determinant divisor. This is the
Hilbert--P\'olya/Connes--Deninger obstruction.
\end{enumerate}

The three hypotheses are independent. The Wilson-surface construction
(\Cref{v4-arith:w9-r2:def:F-infty-WS}), the conjugacy to
\(L_{0}^{\mathrm{Ar}}\) (\Cref{v4-arith:w9-r2:thm:conjugacy}), the
spectral identification with \(\xi\)-zeros
(\Cref{v4-arith:w9-r2:thm:spectral-identification}), and the purity
equivalence (\Cref{v4-arith:w9-r2:thm:purity-ATP-RH}) are each
conditional on a stated subset of (1)--(3).

\section{Bibliography}
\label{v4-arith:w9-r2:bibliography}

Primary references.

\begin{description}
\item[Arakelov 1974.] Arakelov, S.~Yu. \emph{Intersection theory of
divisors on an arithmetic surface}. Izv.~Akad.~Nauk SSSR
Ser.~Mat.~\textbf{38} (1974), 1179--1192.
\item[Baez--Schreiber 2005.] Baez, J.~C.; Schreiber, U.
\emph{Higher gauge theory}. arXiv:hep-th/0412325 (2005).
\item[Berry--Keating 1999.] Berry, M.~V.; Keating, J.~P.
\emph{The Riemann zeros and eigenvalue asymptotics}.
SIAM Rev.~\textbf{41} (1999), 236--266.
\item[Bombieri 2000.] Bombieri, E. \emph{Problems of the millennium:
the Riemann hypothesis}. Clay Mathematics Institute Millennium Prize
Problem statement (2000); revised 2004.
\item[Borger 2009.] Borger, J. \emph{Lambda-rings and the field with
one element}. arXiv:0906.3146 (2009).
\item[Cattaneo--Mnev--Reshetikhin 2014.] Cattaneo, A.~S.; Mnev, P.;
Reshetikhin, N. \emph{Classical BV theories on manifolds with
boundary}. Commun.~Math.~Phys.~\textbf{332} (2014), 535--603.
\item[Connes 1999.] Connes, A. \emph{Trace formula in noncommutative
geometry and the zeros of the Riemann zeta function}. Selecta
Math.~(N.S.)~\textbf{5} (1999), 29--106.
\item[Connes--Marcolli 2008.] Connes, A.; Marcolli, M.
\emph{Noncommutative geometry, quantum fields and motives}. AMS
Colloquium Publications \textbf{55} (2008).
\item[Deninger 1998.] Deninger, C. \emph{Some analogies between
number theory and dynamical systems on foliated spaces}.
Proc.~ICM~Berlin \textbf{1} (1998), 163--186.
\item[Elitzur--Moore--Schwimmer--Seiberg 1989.] Elitzur, S.; Moore,
G.; Schwimmer, A.; Seiberg, N. \emph{Remarks on the canonical
quantization of the Chern-Simons-Witten theory}.
Nucl.~Phys.~B~\textbf{326} (1989), 108--134.
\item[Faltings 1984.] Faltings, G. \emph{Calculus on arithmetic
surfaces}. Ann.~Math.~(2)~\textbf{119} (1984), 387--424.
\item[Hadamard 1893.] Hadamard, J. \emph{Etude sur les propri\'et\'es
des fonctions enti\`eres et en particulier d'une fonction consid\'er\'ee
par Riemann}. J.~Math.~Pures Appl.~(4)~\textbf{9} (1893), 171--215.
\item[Kac--Raina 1987.] Kac, V.~G.; Raina, A.~K. \emph{Bombay
Lectures on Highest Weight Representations of Infinite Dimensional
Lie Algebras}. World Scientific, 1987.
\item[Kim 2015.] Kim, M. \emph{Arithmetic Chern-Simons theory I}.
arXiv:1510.05818 (2015).
\item[Morishita 2012.] Morishita, M. \emph{Knots and Primes: An
Introduction to Arithmetic Topology}. Universitext, Springer, 2012;
arXiv:0904.3399.
\item[Odlyzko 1987.] Odlyzko, A.~M. \emph{On the distribution of
spacings between zeros of the zeta function}.
Math.~Comp.~\textbf{48} (1987), 273--308.
\item[Quillen 1985.] Quillen, D. \emph{Determinants of
Cauchy-Riemann operators over a Riemann surface}.
Funct.~Anal.~Appl.~\textbf{19} (1985), 31--34.
\item[Schreiber--Sati 2011.] Schreiber, U.; Sati, H.
\emph{Survey of mathematical foundations of QFT and perturbative
string theory}. arXiv:1109.0955 (2011); higher-form construction
referenced via arXiv:1011.4735.
\item[Segal 1981.] Segal, G. \emph{Unitary representations of some
infinite dimensional groups}.
Commun.~Math.~Phys.~\textbf{80} (1981), 301--342.
\item[Tate 1950.] Tate, J.~T. \emph{Fourier analysis in number
fields, and Hecke's zeta-functions}. Princeton Ph.D.~thesis, 1950;
reprinted in Cassels--Fr\"ohlich, \emph{Algebraic Number Theory},
Academic Press, 1967.
\item[Weil 1952.] Weil, A. \emph{Sur les ``formules explicites'' de
la th\'eorie des nombres premiers}.
Bull.~Soc.~Math.~France~\textbf{80} (1952), 183--192.
\item[Witten 1989.] Witten, E. \emph{Quantum field theory and the
Jones polynomial}. Commun.~Math.~Phys.~\textbf{121} (1989),
351--399.
\end{description}

\section*{Archimedean Surface Boundary}
\label{v4-arith:w9-r2:summary}

\begin{enumerate}
\item The archimedean Frobenius \(F_{\infty}\) admits an unconditional
Arakelov presentation and two conditional presentations
on the archimedean Fock factor
\(\mathcal{H}^{(\infty)}=\mathcal{V}^{\mathrm{prim}}_{\mathrm{Heis},\infty}\):
the Arakelov scaling exponential \(F_{\infty}=e^{-\Theta_{\infty}}\)
(\Cref{v4-arith:w9-r2:def:F-infty}), the bulk Wilson-surface
observable \(F_{\infty}^{\mathrm{WS}}\) on
\(\Sigma_{\infty}=\{\infty\}\times D^{2}\) with area-scaling
2-representation \(\sigma_{\mathrm{area}}\)
(\Cref{v4-arith:w9-r2:def:F-infty-WS}), and the Mellin-conjugate of
the arithmetic Virasoro zero-mode \(L_{0}^{\mathrm{Ar}}\) under the
Weil--Petersson isometry \(\Upsilon_{\infty}\)
(\Cref{v4-arith:w9-r2:thm:conjugacy}).
\item The three presentations coincide only under the stated
conditional hypotheses:
\(F_{\infty}^{\mathrm{WS}}=F_{\infty}\) by
\Cref{v4-arith:w9-r2:thm:WS-equals-scaling},
\(F_{\infty}=\Upsilon_{\infty}^{-1}e^{-L_{0}^{\mathrm{Ar}}}\Upsilon_{\infty}\)
by \Cref{v4-arith:w9-r2:thm:conjugacy}.
\item Under H-P\(^{\mathrm{Ar}}\), the regularised determinant divisor of
\(F_{\infty}\) on \(H^{1}_{\mathrm{ch}}\) is
\((e^{-\bullet})_{*}\operatorname{div}\xi_{\mathbb R}\)
(\Cref{v4-arith:w9-r2:thm:spectral-identification}). The Hadamard
product belongs to the additive determinant
\(\det_{\infty}(s\mathrm{id}-\Theta_{\xi})\), not to the
multiplicative determinant
\(\det_{\infty}(1-e^{-s}F_{\infty})\)
(\Cref{v4-arith:w9-r2:cor:character-factorisation}).
\item The archimedean \(F_{\infty}\)-weight \(=1/2\) is equivalent to
Deninger archimedean purity, to A-TP, and to RH
(\Cref{v4-arith:w9-r2:thm:purity-ATP-RH}). The concentration of the
Riemann Hypothesis in the archimedean fibre is thereby explained:
finite-place \(\psi_{p}\) components are already pure by Bost--Connes
structure, and archimedean purity of \(F_{\infty}\) is the residual
weight-\(1/2\) condition.
\item The residual obstructions
\Cref{v4-arith:w8-a:rem:wilson-infinity} (Wilson-surface
interpretation of \(F_{\infty}\)) and
\Cref{v4-arith:w8-a:lem:L0-conjugacy} (conjugacy
\(F_{\infty}\sim e^{-L_{0}^{\mathrm{Ar}}}\)) are carried forward to three named
conditionals: quantum ACS bulk measure, primary-Virasoro reconciliation,
the Mellin realisation of \(L_{0}^{\mathrm{Ar}}\), and
H-P\(^{\mathrm{Ar}}\).
\end{enumerate}

The archimedean Frobenius \(F_{\infty}\) as a Wilson-surface operator in 3D
arithmetic Chern-Simons, as the Arakelov scaling exponential, and as the
\(\Upsilon_{\infty}\)-conjugate of \(e^{-L_{0}^{\mathrm{Ar}}}\) become three
descriptions of a single operator only after these residual hypotheses
are imposed. The equivalence of \(F_{\infty}\)-purity
with A-TP locates the Riemann Hypothesis at the archimedean purity
obstruction: finite-place components are pure by the Bost--Connes Adams
structure, and the archimedean component reduces exactly to
Hilbert--P\'olya/Connes--Deninger.
